%قبل از شروع مسئله ترجیحاً چیزی ننویسید. اما اگر خواستید بنویسید، در قالب لیستی (فقرات و شمارش) بنویسید.
%برای اضافه کردن نمره، بعدی از شروع‌مسئله، [x] را اضافه کنید و به جای x نمره مورد نظر را قرار دهید. برای زیرمسئله‌ها هم به همین صورت است.
%\شروع‌مسئله[10]

\شروع‌مسئله [20]
%\متن‌سیاه{سوال عملی اول:}\\
%\متن‌سیاه{دراپ‌تست}\\

% صورت اصلی مسئله را در این‌جا بنویسید.
شنگدباو\پانویس{shengdebao} نواری دارد که شامل $n$ خانه در یک ردیف است و هر خانه را می‌خواهد به رنگی خاص رنگ کند. شنگدباو می‌خواهد این رنگ‌آمیزی را به گونه‌ای انجام دهد که در بازه‌هایی مشخص رنگ‌های تکراری ظاهر نشده باشد. او $q$ بازه با این شرط مشخص کرده است و می‌خواهد \textbf{از کمترین تعداد رنگ} برای این‌کار استفاده کند.

سؤالی که ذهن شنگدباو را مشغول کرده این است که اولاً کمترین تعداد رنگ مختلف برای این‌کار چند‌تاست و چند روش رنگ‌آمیزی مختلف برای این‌کار وجود دارد. از آن‌جایی‌که این تعداد ممکن است خیلی زیاد باشد او می‌خواهد این عدد را به‌پیمانۀ $10^9 + 7$ پیدا کند. 

شنگدباو در مسائل شمارشی اخیراً ناتوان شده است بنابراین از شما می‌خواهد این تعداد را برایش حساب کنید.

\متن‌سیاه{ورودی:}\\
%در خط اول ورودی، ابتدا k و سپس m داده می‌شود.\\
در خط اول ورودی عدد $n$ و $q$ ظاهر شده است که اولی برابر اندازۀ نوار و دومی برابر تعداد بازه‌های متنوع است.

سپس در $q$ خط بعدی در هر خط دو عدد $l$ و $r$ ظاهر شده است که بیانگر این است که بازۀ $[l, r]$ باید رنگ‌هایش متنوع باشند.

\متن‌سیاه{خروجی:}\\
%در تنها سطر خروجی، تعداد آیفون‌های مورد نیاز می‌آید. اگر در هیچ صورتی چنین کاری ممکن نبود ۱- چاپ می‌شود.\\
در تنها خط خروجی ابتدا کمترین تعداد رنگ‌های لازم و کافی برای برقراری شرط‌ها را خروجی دهید و سپس با یک فاصله تعداد روش‌های مختلف رنگ‌آمیزی را به‌پیمانۀ $10^9 + 7$ خروجی دهید.


\متن‌سیاه{محدودیت‌ها:}\\
$1\leq n, q\leq 10^5$

%\textbf{\متن‌سیاه{ورودی و خروجی نمونه}\\

\متن‌سیاه{ورودی و خروجی نمونه:}\\


%% تست اول
\begin{latin}
\begin{tabularx}{\textwidth}{|X | X|}
\hline
\textbf{stdin} &
\textbf{stdout}\\
\hline
3 2

1 2

2 3&
2 2\\
\hline

\end{tabularx}
\end{latin}


%% تست دوم

\begin{latin}
\begin{tabularx}{\textwidth}{|X | X|}
\hline
\textbf{stdin} &
\textbf{stdout}\\
\hline
6 3

1 2

2 3

3 6&
4 216\\
\hline

\end{tabularx}
\end{latin}

%اگر مسئله چند قسمت داشت به این صورت عمل کنید:
%\شروع{زیرمسائل}
%\زیرمسئله[0] % مسئلۀ ۱آ
%\زیرمسئله % مسئلۀ ۱ب
%\زیرمسئله % مسئلۀ ۱پ
%\زیرمسئله % مسئلۀ ۱ت
%\پایان{زیرمسائل}

\solution{
اولاً تعداد رنگ‌های لازم و کافی برابر اندازۀ بزرگترین بازه می‌باشد؛ فرض کنید از ابتدا شروع به رنگ‌کردن می‌کنید و هر سری اولین رنگ ممکن را استفاده می‌کنید. اگر این تعداد رنگ برابر $k$ باشد
هر سری که رنگی جدید می‌خواهیم استفاده کنیم می‌دانیم در بازه‌های شامل این خانه حداکثر $k-1$
خانۀ رنگ‌شده وجود دارد. بنابراین با استفاده از یکی از رنگ‌های دیگر می‌توانیم چنین خانه‌ای را 
رنگ کنیم.

حال سراغ تعداد روش‌ها می‌رویم. فرض کنید همان روند پیشین را انجام می‌دهیم و به ترتیب از خانۀ اول تا $n$ام را رنگ می‌کنیم.
در هر مرحله خانه‌ای را که می‌خواهیم رنگ کنیم در نظر بگیرید. در بین تمامی بازه‌های شامل این خانه آن بازه را در نظر بگیرید که خانه‌های رنگ شده‌اش بیشینه است یعنی \textbf{ابتدای بازه مینیمم است}.
در این‌صورت اگر اندازۀ این بازه را $s$
در نظر بگیرید. تعداد روش‌های رنگ کردن این خانه $k-s$ است.

بنابراین اگر به ازای هر خانه عددی نسبت می‌دهیم. اگر این عدد در هیچ بازه‌ای ظاهر نشده‌بود عدد متناظرش 
$k$
و در غیر این‌صورت 
فرض کنید $m$ برابر مینیمم ابتدا در بین بازه‌هایی باشد
که خانۀ مورد نظر در آن ظاهر شده باشد و خانۀ مورد نظر $i$امین خانه باشد،
عدد متناظرش برابر $k - (i - m)$
خواهد بود.

اگر اعداد متناظر با هر خانه را در هم ضرب کنیم و به‌پیمانۀ عدد مورد نظر حساب کنیم در واقع تعداد حالات را محاسبه کرده‌ایم چرا که عدد متناظر هر خانه برابر تعداد انتخاب‌های رنگی مختلف برای آن خانه‌ است و چون این رنگ‌ها مستقل از هم رفتار می‌کنند با اصل ضرب آن‌ها را می‌توان در هم ضرب کرد.

با داشتن $k$
و به‌دست آوردن $m$
به ازای هر خانه می‌توانیم جواب را محاسبه کنیم.

حال برای محاسبۀ $m$به ازای هر خانه می‌توانیم
از خانۀ آخر به اول حرکت کنیم و هر مرحله اگر به انتهای یکی از بازه‌ها رسیدیم $m$ را برابر مینیمم خودش و ابتدای بازۀ مورد نظر کنیم.
با کمی تفکر در می‌یابید با این‌کار $m$ در هر مرحله عدد مورد نیاز خواهد شد.
و با ضرب $k - (i - m)$
به ازای هر خانه در هم جواب به‌دست می‌آید.

}




%%%%%%%%%%%%%%%%%%%%%%%%%%%%%%%%%%%%%%%%%%%%%%%%%%%%%
% برخی دستورات مفید برای تمرین‌ها. %
%%%%%%%%%%%%%%%%%%%%%%%%%%%%%%%%%%%%%%%%%%%%%%%%%%%%%

% فرمول‌های نمونه:
%$f(x) = \Theta(x)$
%$T(x, y) \leq \log(x) + 2^y + \binom{2n}{n}$
% {\tt code\_function}


% لیست موارد بدون شماره‌گذاری:
%\شروع{فقرات}
% \فقره مورد اول 
% \شروع{فقرات}
% \فقره مورد اول داخلی 
% \شروع{فقرات}
% \فقره اولین مورد خیلی داخلی 
% \فقره دومین مورد خیلی داخلی 
% \پایان{فقرات}
% \فقره مورد دوم داخلی 
% \پایان{فقرات}
% \فقره مورد دوم 
%\پایان{فقرات}

% لیست موارد شماره‌گذاری شده:
%\شروع{شمارش}
% \فقره یک مورد 
% \فقره یک مورد دیگر 
% \فقره باز هم یک مورد دیگر
%\پایان{شمارش}

% برای نوشتن یک فرمول چند خطی که یک شماره بخوره. برای حالت‌های دیگر به سایت زیر مراجعه کنید:
% https://www.overleaf.com/learn/latex/Aligning_equations_with_amsmath
%\begin{align} 
% \begin{split}
% (x+y)^3 &= (x+y)^2(x+y) \\
% &= (x^2+2xy+y^2)(x+y) \\
% &= (x^3+2x^2y+xy^2) + (x^2y+2xy^2+y^3) \\
% &= x^3+3x^2y+3xy^2+y^3
% \end{split} 
%\end{align}

% نوشتن ماتریس و آرایه دوبعدی
%\begin{align}
% A = 
% \begin{bmatrix}
% A_{11} & A_{21} \\
% A_{21} & A_{22}
% \end{bmatrix}
%\end{align}

% اضافه کردن تصویر با شرح. برای اضافه کردن تصویر بدون شرح، فقط دستور درج‌تصویر را استفاده کنید.
%\شروع{شکل}[h!]
%\تنظیم‌ازوسط
% \درج‌تصویر[پهنا=0.5\پهنای‌متن]{img.jpg}
%\شرح{شرح تصویر}
%\برچسب{تصویر:۱}
%\پایان{شکل}
%\پایان{زیرمسائل}

%کد زیر را ببینید.
%%\lstinputlisting{source_filename.py}
%\begin{latin}
%\begin{lstlisting}
%# Import the modules
%import sys
%import random
%
%ans = True
%
%while ans:
%question = raw_input("Ask the magic 8 ball a question: (press enter to quit) ")
%
%answers = random.randint(1,8)
%
%if question == "":
%sys.exit()
%
%elif answers == 1:
%print "It is certain"
%
%elif answers == 2:
%print "Outlook good"
%
%elif answers == 3:
%print "You may rely on it"
%
%elif answers == 4:
%print "Ask again later"
%
%elif answers == 5:
%print "Concentrate and ask again"
%
%elif answers == 6:
%print "Reply hazy, try again"
%
%elif answers == 7:
%print "My reply is no"
%
%elif answers == 8:
%print "My sources say no"
%\end{lstlisting}
%\end{latin}